\documentclass[12pt,a4paper]{article}
\usepackage{fontspec}
\setmainfont{Times New Roman}
\usepackage[top=2cm,bottom=2cm,left=2.5cm,right=2cm]{geometry}
\usepackage{enumitem}
\usepackage{tabularx}
\usepackage{booktabs}
\usepackage{array}
\newcolumntype{Y}{>{\centering\arraybackslash}X}
\pagestyle{plain}
\linespread{1.15}

\begin{document}
% --- HEADER 2 CỘT CHUẨN NGHỊ ĐỊNH 30/2020/NĐ-CP ---
\noindent
\begin{minipage}[t]{0.45\textwidth}
    \begin{center}
        \fontsize{11pt}{13pt}\selectfont
        CÔNG AN THÀNH PHỐ HÀ NỘI\\
        \textbf{CƠ QUAN CẢNH SÁT ĐIỀU TRA (PC02)}\\[-0.2em]
        \rule{3.2cm}{0.6pt}\\[0.25cm]
        \fontsize{11pt}{13pt}\selectfont Số: \textbf{14/BB-KNHT}
    \end{center}
\end{minipage}
\hfill
\begin{minipage}[t]{0.52\textwidth}
    \begin{center}
        \fontsize{11pt}{13pt}\selectfont
        \textbf{CỘNG HÒA XÃ HỘI CHỦ NGHĨA VIỆT NAM}\\
        \textbf{Độc lập - Tự do - Hạnh phúc}\\[-0.2em]
        \rule{3.6cm}{0.6pt}\\[0.25cm]
        \textit{Hà Nội, ngày 25 tháng 07 năm 2016}
    \end{center}
\end{minipage}

\vspace{0.6cm}

\begin{center}
    \textbf{\Large BIÊN BẢN KHÁM NGHIỆM HIỆN TRƯỜNG}\\[0.15cm]
    \textit{}
\end{center}

\vspace{0.4cm}

Vào hồi 07 giờ 15 phút, ngày 25 tháng 07 năm 2016.

Tại: Nhà riêng số 14 Đường Bờ Sông, Phường Phân khu Cảng, Thành phố Hà Nội.



\subsection*{Chúng tôi gồm:}


\textbf{Cán bộ chủ trì khám nghiệm:}

\begin{itemize}[leftmargin=1.2cm, itemsep=2pt]
  \item Ông: \textbf{Lê Văn Dũng} – Điều tra viên thuộc Phòng Cảnh sát Hình sự (PC02) – Công an TP. Hà Nội.

\end{itemize}
\textbf{Cán bộ kỹ thuật hình sự và giám định viên:}

\begin{itemize}[leftmargin=1.2cm, itemsep=2pt]
  \item Ông: \textbf{Trần Tuấn Anh} – Giám định viên Phòng Kỹ thuật Hình sự (PC09) – Công an TP. Hà Nội.
  \item Ông: \textbf{Phạm Minh Hoàng} – Kỹ thuật viên khám nghiệm hiện trường, PC09.

\end{itemize}
\textbf{Đại diện Viện Kiểm sát nhân dân tham gia kiểm sát khám nghiệm:}

\begin{itemize}[leftmargin=1.2cm, itemsep=2pt]
  \item Ông: \textbf{Đỗ Quốc Huy} – Kiểm sát viên Viện Kiểm sát nhân dân TP. Hà Nội.

\end{itemize}
\textbf{Người chứng kiến:}

\begin{itemize}[leftmargin=1.2cm, itemsep=2pt]
  \item Bà: \textbf{Nguyễn Thị Lụa} – Cư trú tại: Tổ dân phố Phân khu Cảng, TP. Hà Nội (Người phát hiện và trình báo).

\end{itemize}
\textit{Cùng sự tham gia bảo vệ hiện trường của đại diện Công an Phường Phân khu Cảng.}



Căn cứ Điều 150 Bộ luật Tố tụng hình sự nước Cộng hòa xã hội chủ nghĩa Việt Nam năm 2003, tiến hành khám nghiệm hiện trường vụ việc phát hiện tử thi anh Nguyễn Văn Khang tại địa chỉ trên.




\subsection*{I. MÔ TẢ HIỆN TRƯỜNG CHUNG}


\begin{itemize}[leftmargin=1.2cm, itemsep=2pt]
  \item Hiện trường là căn phòng khách tầng 1 tại nhà riêng số 14 Đường Bờ Sông. Cửa chính mở hướng ra mặt đường Bờ Sông, không có dấu vết cạy phá cửa.
  \item Không gian bên trong phòng khách có dấu hiệu xáo trộn, đồ đạc lộn xộn. Trên sàn nhà phát hiện nhiều vệt máu và dấu vết vật chất phân bố tập trung ở khu vực bàn tiếp khách và vị trí thi thể.


\end{itemize}

\subsection*{II. MÔ TẢ TỬ THI TẠI HIỆN TRƯỜNG}


\begin{itemize}[leftmargin=1.2cm, itemsep=2pt]
  \item Thi thể nạn nhân (được xác định là anh Nguyễn Văn Khang, sinh năm 1988) nằm ngửa trên sàn phòng khách, đầu chếch về hướng góc bàn tiếp khách.
  \item Quanh vùng đầu và cổ có vũng máu đọng sẫm màu đã đông kết.
  \item \textbf{Quan sát sơ bộ tại chỗ:}
  \item Vùng chẩm gáy bên phải có vết rách dập da tương ứng với cạnh bàn;
  \item Vùng cổ trái có vết thương hở sâu đang chảy máu nhiều làm ướt đẫm phần cổ áo;
  \item Phát hiện một mảng tóc mai bên trái bị cắt sát da đầu.


\end{itemize}

\subsection*{III. DẤU VẾT VÀ VẬT CHỨNG THU THẬP ĐƯỢC TẠI HIỆN TRƯỜNG}


Quá trình khám nghiệm đã phát hiện, đánh số ký hiệu và tiến hành thu giữ các đồ vật, dấu vết sau:


\begin{itemize}[leftmargin=1.2cm, itemsep=2pt]
  \item \textbf{Ký hiệu (41):} Dấu vết đường vân tay chưa rõ danh tính, được trích thu bằng bột kỹ thuật trên bề mặt khung ảnh phòng khách.
  \item \textbf{Ký hiệu (16):} Tập tài liệu, đơn khởi kiện liên quan đến vụ án tranh chấp đất đai, thu giữ trên mặt bàn phòng khách.
  \item \textbf{Ký hiệu (39):} Các mảnh vỡ bình trà bằng gốm; trong đó có 01 mảnh thủy tinh nhọn dài dính nhiều máu nằm trên sàn nhà nghi là hung khí gây thương tích vùng cổ.
  \item \textbf{Ký hiệu (40):} 01 khung ảnh để bàn bị rơi xuống sàn, mặt kính vỡ rạn và nứt vỡ nhiều mảnh.
  \item \textbf{Ký hiệu (18):} Nhiều mảnh giấy báo bị xé vụn nằm vương vãi trên nền nhà phía chân bàn.
  \item \textbf{Ký hiệu (13):} 01 cuốn sổ tay kích thước nhỏ, bên trong ghi chép các khoản nợ tiền.
  \item \textbf{Ký hiệu (dev-00):} 01 điện thoại di động nhãn hiệu Apple iPhone 6s Plus của nạn nhân, thu giữ cạnh vị trí cánh tay nạn nhân.

\end{itemize}
Toàn bộ đồ vật, dấu vết nêu trên đã được mô tả vào sơ đồ, chụp ảnh hiện trường và niêm phong theo đúng quy chuẩn.



Biên bản khám nghiệm kết thúc hồi 08 giờ 30 phút cùng ngày, đã được đọc lại cho mọi người có tên trên cùng nghe, công nhận đúng và ký tên dưới đây.




\vspace{0.8cm}
\noindent
\begin{tabularx}{\textwidth}{YYYY}
\toprule
\textbf{CÁN BỘ CHỦ TRÌ KHÁM NGHIỆM} & \textbf{CÁN BỘ KỸ THUẬT HÌNH SỰ} & \textbf{KIỂM SÁT VIÊN KIỂM SÁT} & \textbf{NGƯỜI CHỨNG KIẾN} \\
\midrule
\textit{(Ký, ghi rõ họ tên)} & \textit{(Ký, ghi rõ họ tên)} & \textit{(Ký, ghi rõ họ tên)} & \textit{(Ký, ghi rõ họ tên)} \\[1.5cm]
\textbf{Lê Văn Dũng} & \textbf{Trần Tuấn Anh} & \textbf{Đỗ Quốc Huy} & \textbf{Nguyễn Thị Lụa} \\
\bottomrule
\end{tabularx}
\end{document}
